\subsection{Data Sources and Search Protocol}
\label{sec:search-protocol}
We searched five electronic sources: OpenAlex\cite{b14}, SpringerLink\cite{b15}, IEEE Xplore\cite{b18}, arXiv\cite{b16}, and PubMed\cite{b17}. The searches were initiated on May 1, 2026. The review considered English-language records published or posted from January 1, 2021, through the evidence cutoff of May 31, 2026. The search was designed around three concept groups: (i) retrieval-augmented generation or LLM-based generation, (ii) knowledge graphs or semantic graphs, and (iii) ontology or ontological constraints. This conjunction was intended to retrieve both ontology-aware KG-RAG candidates and graph-linked ontology-grounded LLM systems adjacent to RAG. The source totals sum to 1,367 records: OpenAlex (576), SpringerLink (594), IEEE Xplore (33), arXiv (151), and PubMed (13). A consolidated record-level search log is retained in the analysis package for all 1,367 captured records. The log records source designation, title, authors, publication information, abstract or snippet, and DOI where available, thereby supporting inspection of the captured set and source totals. Because the log does not include original provider exports, interface snapshots, or complete page/batch logs, exact replication of query execution and the stopping process remains limited.

OpenAlex was used as a broad open bibliographic discovery layer because it aggregates scholarly metadata from multiple sources, including Crossref, PubMed, arXiv, and other repositories. This choice improves reproducibility because OpenAlex provides open metadata access and does not require a proprietary subscription. However, OpenAlex should not be interpreted as a complete substitute for curated subject or citation databases such as Scopus, or for publisher-specific full-text collections such as the ACM Digital Library. Because ACM Digital Library and Scopus were not searched as separate sources, relevant computing, semantic-web, or information-retrieval papers indexed only or more reliably in those databases may have been missed. We therefore treat this as a search-coverage limitation rather than claiming exhaustive database coverage.

The master Boolean expression was adapted to the syntax and export options of each source. The database-specific search strings were as follows:

\begin{enumerate}
    \item \textbf{OpenAlex} (metadata search; $n=576$): \texttt{(ontology OR ontological OR "ontology-aware" OR "ontology-driven") AND ("knowledge graph" OR KG OR "semantic graph") AND ("retrieval-augmented generation" OR RAG OR "knowledge-augmented generation" OR "large language model")}


    \item \textbf{SpringerLink} (advanced search; $n=594$): \texttt{(ontology OR ontological OR "ontology-aware" OR "ontology-driven") AND ("knowledge graph" OR KG OR "semantic graph") AND ("retrieval-augmented generation" OR RAG OR "knowledge-augmented generation" OR "large language model")}

    \item \textbf{IEEE Xplore} (advanced search; $n=33$): \texttt{(ontology OR ontological OR "ontology-aware" OR "ontology-driven") AND ("knowledge graph" OR KG OR "semantic graph") AND ("retrieval-augmented generation" OR RAG OR "knowledge-augmented generation" OR "large language model")}


    \item \textbf{arXiv} (advanced search in title/abstract/all fields; $n=151$): \texttt{(ontology OR ontological OR "ontology-aware" OR "ontology-driven") AND ("knowledge graph" OR KG OR "semantic graph") AND ("retrieval-augmented generation" OR RAG OR "knowledge-augmented generation" OR "large language model")}

    \item \textbf{PubMed} (title/abstract search; $n=13$): \texttt{(ontology[Title/Abstract] OR ontological[Title/Abstract] OR ontology-aware[Title/Abstract] OR ontology-driven[Title/Abstract]) AND ("knowledge graph"[Title/Abstract] OR "semantic graph"[Title/Abstract] OR "linked data"[Title/Abstract]) AND ("retrieval-augmented generation"[Title/Abstract] OR "retrieval augmented generation"[Title/Abstract] OR RAG[Title/Abstract] OR "knowledge-augmented generation"[Title/Abstract] OR "large language model"[Title/Abstract])}
\end{enumerate}

We examined search-result pages from sources with large result sets sequentially under the stated date-range, language, and content-type filters. The operational stopping rule was to stop when an additional result page contained no records matching all three concept groups with usable title, author, and year metadata. No automated scraping was used. Records were entered using a common metadata schema covering visible title, author, year, venue, URL, and abstract or snippet; DOI or arXiv identifiers were added when visible in the record or full text. Where bulk export was unavailable or incomplete, records were entered manually from the source result pages using the same fields. The consolidated log permits record-level inspection of the retained search set, although the stopping procedure cannot be replayed without complete page-level logs.

OpenAlex records were obtained from metadata search results; SpringerLink and IEEE Xplore records were exported or checked manually from search-result metadata pages where bulk export was limited; and arXiv and PubMed records were obtained from citation or metadata pages. The record-level log supports verification of the captured metadata and source counts. Original query exports, interface snapshots, and page-level stopping logs would still be required to reproduce the source interfaces and search execution exactly.